% Sheaves and Real-World Information: an introduction and reading guide -- LuaLaTeX.  Build: sh build.sh
\documentclass[11pt,a4paper,oneside]{scrartcl}

% ---- page and type ------------------------------------------------------------
\usepackage[textwidth=31pc,textheight=47\baselineskip,top=1.1in,headsep=2.2\baselineskip,
            footskip=3\baselineskip,heightrounded]{geometry}
\usepackage{amsmath,amsthm,mathtools}
\usepackage{unicode-math}
\setmainfont{Libertinus Serif}[Numbers=OldStyle]
\setsansfont{Libertinus Sans}
\setmathfont{Libertinus Math}
\setmathfont{Latin Modern Math}[range={\mathcal,\mathbfcal,\mathscr}]
\usepackage[babel=true]{microtype}
\usepackage[english]{babel}
\usepackage[autostyle]{csquotes}
\usepackage{booktabs,tabularx,array,enumitem}
\usepackage[backend=biber,style=authoryear,maxcitenames=2,maxbibnames=10,giveninits=true,
            uniquename=false,doi=false,url=false,isbn=false,dashed=false]{biblatex}
\addbibresource{refs.bib}
\renewcommand*{\bibfont}{\small}
\setlength{\bibhang}{1.4em}
\renewbibmacro{in:}{\ifentrytype{article}{}{\printtext{\bibstring{in}\intitlepunct}}}
\usepackage{scrlayer-scrpage}
\usepackage{flafter}
\usepackage{tikz}
\usetikzlibrary{arrows.meta,calc}
\usepackage[hidelinks,pdfusetitle]{hyperref}

\linespread{1.05}
\setlength{\parindent}{1.4em}
\widowpenalty=4000 \clubpenalty=4000
\setlength{\emergencystretch}{1.5em}
\setlist{itemsep=.25em,topsep=.5em,parsep=0pt}

% ---- running heads --------------------------------------------------------------
\clearpairofpagestyles
\ohead{\textsc{\MakeLowercase{Sheaves and Real-World Information}}}
\ihead{\textsc{\MakeLowercase{Introduction and Reading Guide}}}
\cfoot*{\pagemark}
\setkomafont{pageheadfoot}{\normalfont\footnotesize}
\setkomafont{pagenumber}{\normalfont\small}
\pagestyle{scrheadings}

% ---- headings and captions ----------------------------------------------------
\setkomafont{disposition}{\normalfont}
\setkomafont{section}{\normalfont\large\scshape}
\setkomafont{subsection}{\normalfont\normalsize\itshape}
\RedeclareSectionCommand[beforeskip=-2\baselineskip plus -.4\baselineskip,afterskip=.8\baselineskip]{section}
\RedeclareSectionCommand[beforeskip=-1.2\baselineskip plus -.3\baselineskip,afterskip=.4\baselineskip]{subsection}
\renewcommand*{\sectionformat}{\thesection.\enskip}
\renewcommand*{\subsectionformat}{\thesubsection.\enskip}
\setkomafont{caption}{\small}
\setkomafont{captionlabel}{\scshape}
\setcapindent{0pt}
\renewcommand*{\captionformat}{.\enskip}

% ---- statements (given without proofs) ------------------------------------------
\newtheoremstyle{guideplain}{.8\baselineskip}{.8\baselineskip}{\itshape}{}{\scshape}{.}{.5em}{}
\newtheoremstyle{guidedef}{.8\baselineskip}{.8\baselineskip}{\upshape}{}{\scshape}{.}{.5em}{}
\theoremstyle{guideplain}
\newtheorem{theorem}{Theorem}[section]
\newtheorem{proposition}[theorem]{Proposition}
\theoremstyle{guidedef}
\newtheorem{definition}[theorem]{Definition}
\newtheorem{example}[theorem]{Example}
\newtheorem{remark}[theorem]{Remark}

% ---- notation ----------------------------------------------------------------------
\newcommand{\F}{\mathcal{F}}
\newcommand{\G}{\mathcal{G}}
\newcommand{\U}{\mathcal{U}}
\newcommand{\R}{\mathbb{R}}
\DeclareMathOperator{\Op}{Op}
\DeclareMathOperator{\im}{im}
\newcommand{\face}{\trianglelefteq}

% ---- figures -----------------------------------------------------------------------
\colorlet{soft}{black!45}
\tikzset{
  dot/.style  = {circle, fill, inner sep=0pt, minimum size=3.6pt},
  thin/.style = {line width=.5pt},
  arr/.style  = {-{Stealth[length=2mm,width=1.6mm]}, line width=.5pt},
  lab/.style  = {font=\footnotesize},
  soft/.style = {font=\footnotesize, text=soft},
}

\hypersetup{pdftitle={Sheaves and Real-World Information: An Introduction and Reading Guide},
            pdfauthor={Mario Asis}}

% ===================================================================================
\begin{document}

\thispagestyle{plain.scrheadings}
\begin{center}
  \vspace*{.2\baselineskip}
  {\LARGE Sheaves and Real-World Information\par}
  \vspace{.7\baselineskip}
  {\large\itshape An Introduction and Reading Guide\par}
  \vspace{1.4\baselineskip}
  {\scshape \MakeLowercase{Mario Asis}\par}
  \vspace{.3\baselineskip}
  {\small October 2026\par}
  \vspace{1.6\baselineskip}
\end{center}

\begin{center}
\begin{minipage}{26pc}
  \small
  \noindent{\scshape summary.}\enspace
  Sheaves were introduced in the 1940s to pass from local to global information in topology, and
  for sixty years they belonged almost entirely to pure mathematics. Since about 2010 a finite,
  computable version of the theory has been used to integrate measurements from real sensors, to
  describe entropy and the consistency of probabilistic data, and to build neural networks. These
  notes recall what a sheaf is and where the idea came from, explain what changed to make sheaves
  usable on actual data, state the main definitions and results without proofs, and end with an
  ordered reading guide to the literature.
\end{minipage}
\end{center}
\vspace{.8\baselineskip}

{\small\setcounter{tocdepth}{1}\renewcommand{\contentsname}{\normalsize\scshape contents}\tableofcontents}

% -----------------------------------------------------------------------------------
\section{The idea in one paragraph}\label{sec:idea}

A sheaf assigns data to the parts of a space and records how data on a larger part restrict to a
smaller one, in such a way that compatible local data determine a unique global datum. Its
cohomology measures the failure of that last step: the local data that cannot be assembled, and the
ambiguity when they can. For a mathematician the standard examples are functions, differential forms,
and solutions of differential equations. The applied theory reads the same structure differently.
The \enquote{space} is a network of sensors, variables, or agents; the data on each part are
measurements, probability distributions, or feature vectors; the restriction maps say how one
observation is to be compared with another; a global section is a globally consistent account of all
the observations; and cohomology, or a metric refinement of it, quantifies inconsistency. Everything
below elaborates this dictionary.

% -----------------------------------------------------------------------------------
\section{What a sheaf is}\label{sec:def}

\subsection{Presheaves and sheaves}

Let $X$ be a topological space and $\Op(X)$ its poset of open sets, ordered by inclusion.

\begin{definition}[Presheaf]
  A \emph{presheaf} of sets on $X$ is a functor $\F\colon \Op(X)^{\mathrm{op}}\to\mathbf{Set}$. It
  assigns a set $\F(U)$ of \emph{sections} to each open $U$, and to each inclusion $V\subseteq U$ a
  \emph{restriction map} $\rho^U_V\colon\F(U)\to\F(V)$, written $s\mapsto s|_V$, such that
  $\rho^U_U=\mathrm{id}$ and $\rho^V_W\circ\rho^U_V=\rho^U_W$ for $W\subseteq V\subseteq U$. Replacing
  $\mathbf{Set}$ by abelian groups, vector spaces, or rings gives presheaves with values in those
  categories.
\end{definition}

\begin{definition}[Sheaf]\label{def:sheaf}
  A presheaf $\F$ is a \emph{sheaf} if for every open $U$ and every open cover $U=\bigcup_i U_i$:
  \begin{enumerate}[label=(\roman*)]
    \item \emph{(locality)} if $s,t\in\F(U)$ satisfy $s|_{U_i}=t|_{U_i}$ for all $i$, then $s=t$;
    \item \emph{(gluing)} if $s_i\in\F(U_i)$ satisfy $s_i|_{U_i\cap U_j}=s_j|_{U_i\cap U_j}$ for all
          $i,j$, then there is $s\in\F(U)$ with $s|_{U_i}=s_i$ for all $i$.
  \end{enumerate}
  Equivalently, the diagram
  $\F(U)\to\prod_i\F(U_i)\rightrightarrows\prod_{i,j}\F(U_i\cap U_j)$ is an equalizer.
\end{definition}

\begin{figure}[tp]
  \centering
  \begin{tikzpicture}[x=1cm,y=1cm]
    % base
    \draw[line width=1.4pt] (0.6,0) -- (5.4,0);
    \draw[line width=1.4pt] (3.6,-0.28) -- (8.4,-0.28);
    \node[lab] at (1.4,0.27) {$U$};
    \node[lab] at (7.6,-0.57) {$V$};
    \draw[soft, densely dotted, thin] (3.6,-0.28) -- (3.6,2.55);
    \draw[soft, densely dotted, thin] (5.4,-0.28) -- (5.4,2.55);
    \node[soft] at (4.5,-0.75) {$U\cap V$};
    % sections
    \draw[line width=.8pt] plot[smooth] coordinates {(0.6,1.2) (1.6,1.9) (2.8,1.5) (3.6,1.7) (4.5,2.2) (5.4,1.9)};
    \draw[line width=.8pt, dashed] plot[smooth] coordinates {(3.6,1.7) (4.5,2.2) (5.4,1.9) (6.4,1.3) (7.4,1.6) (8.4,1.1)};
    \node[lab] at (1.5,2.25) {$s\in\F(U)$};
    \node[lab] at (7.5,2.0) {$t\in\F(V)$};
    \node[lab] at (4.5,2.85) {$s|_{U\cap V}=t|_{U\cap V}$};
  \end{tikzpicture}
  \caption{Gluing. Two sections that agree on the overlap determine exactly one section over
    $U\cup V$. Locality gives uniqueness, gluing gives existence.}
  \label{fig:gluing}
\end{figure}

\begin{example}
  Continuous real functions, $U\mapsto C(U,\R)$, form a sheaf, as do smooth functions, holomorphic
  functions, sections of a vector bundle, and solutions of a linear differential equation. Bounded
  functions form a presheaf that is not a sheaf: bounded functions on the members of a cover glue to
  a function that need not be bounded. Constant functions likewise fail to glue on a disconnected
  set, whereas \emph{locally} constant functions form the \emph{constant sheaf} $\underline{A}$.
\end{example}

\begin{definition}[Stalks]
  The \emph{stalk} of $\F$ at $x\in X$ is $\F_x=\varinjlim_{U\ni x}\F(U)$; its elements are
  \emph{germs} of sections at $x$. A morphism of sheaves is an isomorphism if and only if it induces
  isomorphisms on all stalks.
\end{definition}

\begin{proposition}[Sheafification]
  The inclusion of sheaves into presheaves has a left adjoint $\F\mapsto\F^{+}$, which preserves
  stalks. Concretely, $\F^{+}(U)$ consists of the families of germs $(s_x)_{x\in U}$ that are locally
  induced by sections of $\F$.
\end{proposition}

\subsection{Cohomology}

For a sheaf of abelian groups the global-sections functor $\Gamma(X,-)\colon\F\mapsto\F(X)$ is left
exact but not exact: a surjection of sheaves need not be surjective on global sections.

\begin{definition}[Sheaf cohomology]
  $H^i(X,\F)$ is the $i$-th right derived functor of $\Gamma(X,-)$ applied to $\F$. Thus
  $H^0(X,\F)=\F(X)$, and a short exact sequence of sheaves induces a long exact sequence in
  cohomology.
\end{definition}

The working interpretation is that $H^0$ is the space of global solutions and $H^1$ classifies the
obstructions to assembling local solutions into a global one. In practice cohomology is computed
from a cover. For an open cover $\U=\{U_i\}$, the Čech complex has
$\check C^k(\U,\F)=\prod_{i_0<\dots<i_k}\F(U_{i_0}\cap\dots\cap U_{i_k})$ with the alternating-sum
differential.

\begin{theorem}[Leray's theorem on acyclic covers]\label{thm:leray}
  If $H^q(U_{i_0}\cap\dots\cap U_{i_k},\F)=0$ for all $q>0$ and all finite intersections of members
  of $\U$, then $\check H^k(\U,\F)\cong H^k(X,\F)$ for all $k$.
\end{theorem}

Theorem~\ref{thm:leray} is the reason the theory can be made finite: with a good cover, cohomology is
the cohomology of an explicit complex indexed by a combinatorial object, the nerve of the cover.
Section~\ref{sec:cellular} takes that combinatorial object as the starting point.

% -----------------------------------------------------------------------------------
\section{Where sheaves came from}\label{sec:history}

Sheaves were not created by Grothendieck, although the association is understandable: he rebuilt the
theory in the form now taught. The history has three steps
\autocite{miller2000,gray1979,houzel1990}.

\emph{Leray.} Jean Leray, an analyst known for his work on the Navier--Stokes equations, spent 1940
to 1945 as a prisoner of war in the officers' camp Oflag XVII-A in Austria. To avoid work that could
serve the war effort he turned to algebraic topology, and there he devised three things at once:
sheaves (\emph{faisceaux}), their cohomology, and spectral sequences. His aim was topological and
quite specific: to relate the cohomology of a space to that of the fibres of a map defined on it,
with fixed-point theorems for equations in analysis as the motivation. The work became known only
after the war, and several more years passed before it was widely understood.

\emph{Cartan, Oka, Serre.} In his Paris seminar of 1948--51 Henri Cartan reformulated Leray's theory
in terms of open sets, with the étalé-space definition due to Michel Lazard, and it became the
natural language for several complex variables: Oka's coherence theorem and Cartan's Theorems A and B
are statements about sheaf cohomology. In 1955 Serre's \emph{Faisceaux algébriques cohérents} carried
the method into algebraic geometry, where there are no small open sets and sheaf cohomology had to do
the work that classical topology could not \autocite{serre1955}.

\emph{Grothendieck.} The 1957 Tôhoku paper placed sheaf cohomology within homological algebra:
abelian categories, enough injectives, and derived functors, so that cohomology is defined on any
space for any sheaf \autocite{grothendieck1957}. Schemes followed, and then, in the early 1960s,
sites and toposes, which free the notion of sheaf from topological spaces altogether. Godement's
book fixed the textbook form of the classical theory \autocite{godement1958}, and the microlocal
theory of Sato and Kashiwara extended it to analysis \autocite{kashiwara1990}.

% -----------------------------------------------------------------------------------
\section{From pure mathematics to data}\label{sec:applied}

It is not accurate to say that sheaves lay unused. From 1950 onwards they were central to algebraic
geometry, complex analysis, and the theory of differential equations. What arrived late was their use
as a model for \emph{data}, and there were two reasons for the delay. The first is computational: a
sheaf on a topological space involves infinitely many open sets and typically infinite-dimensional
spaces of sections, so that nothing can be computed from measurements. The second is that the
people who had data, and the people who knew sheaves, were not the same people.

Three developments, at first unconnected, prepared the applied theory.
\begin{itemize}
  \item In probability, \textcite{vorobev1962} characterized when a consistent family of marginal
        distributions extends to a joint distribution. This is a gluing theorem stated without
        sheaves (Theorem~\ref{thm:vorobev} below).
  \item In computer science, Goguen proposed from the 1970s that systems of interacting objects be
        modelled by sheaves, with behaviour over time as sections \autocite{goguen1992}.
  \item In topology, \textcite{shepard1985} showed that sheaves constructible with respect to a cell
        structure are equivalent to a finite, purely combinatorial object: a \emph{cellular sheaf}.
\end{itemize}
The last is the decisive technical step. A cellular sheaf is a finite diagram of vector spaces and
linear maps, and its cohomology is the cohomology of a finite complex of matrices. Around 2010 this
was taken up in applied topology, principally by Robert Ghrist and his students, and within a few
years the three lines of Section~\ref{sec:uses} appeared: sensor integration \autocite{robinson2017},
contextuality and entropy \autocite{abramsky2011,baudot2015}, and, from 2020, machine learning
\autocite{hansengebhart2020,bodnar2022}. Table~\ref{tab:timeline} gives the chronology.

\begin{table}[tp]
  \centering\small
  \renewcommand{\arraystretch}{1.22}
  \captionabove{A timeline. The upper part is the classical theory; the lower part, its use for data.}
  \label{tab:timeline}
  \begin{tabularx}{\textwidth}{@{}>{\raggedleft\arraybackslash}p{3.6em}!{\color{soft}\vrule width .6pt}>{\raggedright\arraybackslash}X@{}}
    \toprule
    1940--45 & Leray, as a prisoner of war, devises sheaves, sheaf cohomology, and spectral sequences. \\
    1948--51 & Cartan's seminar: open sets, étalé spaces (Lazard), application to several complex variables. \\
    1955     & Serre, \emph{Faisceaux algébriques cohérents}: sheaves enter algebraic geometry. \\
    1957     & Grothendieck, the Tôhoku paper: derived functors; then schemes, sites, and toposes. \\
    1958     & Godement's textbook. \\
    1990     & Kashiwara and Schapira, \emph{Sheaves on Manifolds}: the microlocal theory. \\
    \midrule
    1962     & Vorob'ev: when consistent marginals extend to a joint distribution. \\
    1985     & Shepard's thesis: cellular sheaves. \\
    1992     & Goguen: sheaf semantics for concurrent interacting objects. \\
    2011     & Ghrist and Hiraoka: network coding as sheaf cohomology. Abramsky and Brandenburger:
               contextuality as the absence of a global section. \\
    2014     & Curry's thesis systematizes cellular sheaves and cosheaves. Robinson, \emph{Topological
               Signal Processing}. Ghrist, \emph{Elementary Applied Topology}. \\
    2015     & Baudot and Bennequin: Shannon entropy as a cohomology class. \\
    2017     & Robinson: sheaves as the canonical data structure for sensor integration. \\
    2019     & Hansen and Ghrist: the sheaf Laplacian and its spectral theory. \\
    2020     & Robinson: the consistency radius. Joslyn et al.: sheaf-based tracking on field data.
               Vigneaux: information cohomology in general. Hansen and Gebhart: sheaf neural networks. \\
    2022     & Bodnar et al.: neural sheaf diffusion, the first results on standard benchmarks. \\
    2025     & Ayzenberg et al.: a 117-page survey of sheaves in deep learning. \\
    \bottomrule
  \end{tabularx}
\end{table}

% -----------------------------------------------------------------------------------
\section{Cellular sheaves: the computable version}\label{sec:cellular}

Let $X$ be a regular cell complex (a graph or a simplicial complex will do) and write
$\sigma\face\tau$ when the cell $\sigma$ is a face of $\tau$.

\begin{definition}[Cellular sheaf]\label{def:cellular}
  A \emph{cellular sheaf} $\F$ of vector spaces on $X$ assigns a vector space $\F(\sigma)$, the
  \emph{stalk}, to each cell $\sigma$, and a linear \emph{restriction map}
  $\F_{\sigma\face\tau}\colon\F(\sigma)\to\F(\tau)$ to each incidence $\sigma\face\tau$, such that
  $\F_{\sigma\face\sigma}=\mathrm{id}$ and
  $\F_{\tau\face\upsilon}\circ\F_{\sigma\face\tau}=\F_{\sigma\face\upsilon}$. In other words, $\F$ is
  a functor from the face poset of $X$ to vector spaces.
\end{definition}

The maps run from faces to cofaces, which is the direction of restriction once cells are identified
with their open stars: the star of an edge is contained in the star of each of its endpoints. A
cellular sheaf is equivalent to a sheaf in the sense of Definition~\ref{def:sheaf} on the face poset
with its Alexandrov topology \autocite{curry2014}.

\begin{definition}[Sections and cohomology]
  A \emph{global section} of $\F$ is a choice $x_\sigma\in\F(\sigma)$ for every cell such that
  $\F_{\sigma\face\tau}x_\sigma=x_\tau$ whenever $\sigma\face\tau$. Let
  $C^k(X;\F)=\bigoplus_{\dim\sigma=k}\F(\sigma)$ and define
  $\delta\colon C^k\to C^{k+1}$ by
  $(\delta x)_\tau=\sum_{\sigma\face\tau}[\sigma:\tau]\,\F_{\sigma\face\tau}x_\sigma$,
  where $[\sigma:\tau]=\pm1$ is the incidence number. Then $\delta^2=0$, and
  $H^k(X;\F)=\ker\delta^k/\im\delta^{k-1}$.
\end{definition}

\begin{proposition}
  $H^0(X;\F)=\ker\delta^0$ is the space of global sections. On a graph, with an edge $e$ oriented
  from $u$ to $v$, $(\delta x)_e=\F_{v\face e}x_v-\F_{u\face e}x_u$, so a $0$-cochain is a section
  exactly when the two endpoints of every edge agree after restriction to the edge.
\end{proposition}

\begin{figure}[tp]
  \centering
  \begin{tikzpicture}[x=1cm,y=1cm]
    \foreach \x/\lab/\sgn/\res in {0/{(a) constant sheaf}/{1}/{$H^0=\R,\ H^1=\R$},
                                   6.2/{(b) one map reversed}/{-1}/{$H^0=0,\ H^1=0$}} {
      \begin{scope}[xshift=\x cm]
        \node[dot, label={[lab]left:$\R$}]  (u) at (0,0) {};
        \node[dot, label={[lab]right:$\R$}] (v) at (3.4,0) {};
        \node[dot, label={[lab]above:$\R$}] (w) at (1.7,2.6) {};
        \draw[thin] (u) -- (v) -- (w) -- (u);
        \node[lab, fill=white, inner sep=1.5pt] (euv) at (1.7,0)    {$\R$};
        \node[lab, fill=white, inner sep=1.5pt] (evw) at (2.55,1.3) {$\R$};
        \node[lab, fill=white, inner sep=1.5pt] (ewu) at (0.85,1.3) {$\R$};
        % restriction maps, written at the end of the edge nearest the vertex they start from
        \node[soft] at (0.50,-0.27) {$1$};
        \node[soft] at (2.90,-0.27) {$1$};
        \node[soft] at (3.36,0.62)  {$1$};
        \node[soft] at (2.34,2.18)  {$1$};
        \node[soft] at (1.06,2.18)  {$1$};
        \ifnum\sgn<0 \node[lab] at (-0.06,0.62) {$\mathbf{-1}$};
        \else        \node[soft] at (0.04,0.62) {$1$}; \fi
        \node[lab] at (1.7,-1.0) {\lab};
        \node[soft] at (1.7,-1.5) {\res};
      \end{scope}
    }
  \end{tikzpicture}
  \caption{Two cellular sheaves on a triangle, with stalk $\R$ on every vertex and edge. The
    number at each end of an edge is the restriction map from that vertex to the edge. (a)~All maps are the identity: sections are the
    constant assignments. (b)~Reversing the sign of a single map leaves every local condition
    satisfiable but removes all non-zero global sections: going once round the cycle returns the
    negative of the starting value.}
  \label{fig:triangle}
\end{figure}

Figure~\ref{fig:triangle} shows the smallest example in which local consistency does not yield a
global section. When the stalks carry inner products, the adjoint of $\delta$ is available.

\begin{definition}[Sheaf Laplacian]
  The \emph{sheaf Laplacian} of $\F$ is $L_{\F}=\delta^{*}\delta\colon C^0\to C^0$, and more generally
  $\Delta^k=(\delta^{k})^{*}\delta^{k}+\delta^{k-1}(\delta^{k-1})^{*}$. For the constant sheaf
  $\underline{\R}$ on a graph, $L_{\F}$ is the usual graph Laplacian.
\end{definition}

\begin{theorem}[\cite{hansen2019}]\label{thm:hodge}
  For a cellular sheaf of finite-dimensional inner-product spaces, $\ker\Delta^k\cong H^k(X;\F)$. In
  particular $\ker L_{\F}$ is the space of global sections, and the solution of the heat equation
  $\dot x=-L_{\F}\,x$ converges, as $t\to\infty$, to the orthogonal projection of $x(0)$ onto the space
  of global sections.
\end{theorem}

Theorem~\ref{thm:hodge} turns the search for the consistent account nearest to given data into a
linear dynamical system that each node can run using only its neighbours. Both the fusion algorithms
of Section~\ref{sec:fusion} and the neural networks of Section~\ref{sec:ml} rest on it.

% -----------------------------------------------------------------------------------
\section{Three uses}\label{sec:uses}

\subsection{Integrating measurements}\label{sec:fusion}

\textcite{robinson2017} argues that any reasonable set of axioms for combining heterogeneous sources
leads to a sheaf. The base is a finite topological space, or a cell complex, whose points are the
sensors and their overlaps; the stalk over a sensor is the set in which its readings take values; the
stalk over an overlap is a common set in which the readings of the overlapping sensors can be
compared; and the restriction maps perform the conversion. The actual readings form an
\emph{assignment}: one element of each stalk, with no consistency required.

\begin{definition}[Consistency radius, \cite{robinson2020}]
  Let $\F$ be a sheaf of pseudometric spaces on a finite space and $a$ an assignment, with
  $a(U)\in\F(U)$ for each open $U$. Its \emph{consistency radius} is
  \[
    c(a)=\sup_{V\subseteq U}\,d_{V}\bigl(\rho^U_V\,a(U),\,a(V)\bigr).
  \]
\end{definition}

\begin{proposition}
  When the stalks are metric spaces, $c(a)=0$ if and only if $a$ is a global section. The consistency
  radius is continuous in the assignment, and thresholding it yields a filtration of the base space
  that is stable under perturbation of the data.
\end{proposition}

Data fusion then becomes an optimization problem: find the global section nearest to the assignment,
with the consistency radius as a certified bound on the residual disagreement. Two treatments with
real measurements are worth reading for the modelling as much as for the mathematics: the location of
a hidden radio transmitter from bearings and signal strengths \autocite{robinson2019}, and the
tracking of wild black bears from several kinds of sensor at once \autocite{joslyn2020}.

\subsection{Information theory}\label{sec:info}

\emph{The marginal problem.} Let variables be the vertices of a simplicial complex $K$ whose faces
are the sets of variables observed jointly. Assigning to each face the probability distributions on
its variables, with marginalization as restriction, gives a presheaf. A family of distributions, one
per maximal face, is \emph{consistent} if its members agree on common marginals.

\begin{theorem}[\cite{vorobev1962}]\label{thm:vorobev}
  Every consistent family of distributions on the faces of $K$ is the family of marginals of a single
  joint distribution if and only if $K$ is regular, that is, acyclic in the sense that it can be
  reduced to a point by successively removing faces attached along a single face.
\end{theorem}

When $K$ is not regular there are consistent families with no joint distribution: local sections
with no global section. \textcite{abramsky2011} showed that this is exactly the structure of
non-locality and contextuality in quantum mechanics, and \textcite{abramsky2013} that the same
question arises in database theory as the existence of a universal relation for a set of tables.
Čech cohomology supplies computable witnesses of the obstruction \autocite{abramsky2015}.

\emph{Entropy.} \textcite{baudot2015} define a cohomology theory for an \emph{information
structure}, a category of observables ordered by refinement, with coefficients in functions of
probability distributions on which observables act by conditional expectation.

\begin{theorem}[\cite{baudot2015,vigneaux2020}]
  The $1$-cocycle condition of information cohomology is the chain rule
  $f(XY)=f(X)+X.f(Y)$. Under mild non-degeneracy assumptions on the information structure, the first
  cohomology group is one-dimensional and is generated by Shannon entropy. For the action deformed by
  a parameter $\alpha\neq1$, it is generated by the Tsallis $\alpha$-entropy.
\end{theorem}

Entropy is thus characterized by a topological invariant instead of by a list of axioms. Two related
results complete the picture: information flows in a network code are the degree-zero cohomology of
a sheaf on the network \autocite{ghristhiraoka2011}, and belief propagation in graphical models is a
diffusion whose fixed points are described homologically \autocite{peltre2020}.

\subsection{Machine learning}\label{sec:ml}

A graph neural network diffuses node features with the graph Laplacian, which by the remark above is
the Laplacian of the constant sheaf. Its limit is the constant section, which is the phenomenon known
as oversmoothing, and it implicitly assumes that neighbours should agree, which fails on
\emph{heterophilic} graphs. \textcite{hansengebhart2020} proposed replacing the constant sheaf by a
general cellular sheaf, and \textcite{bodnar2022} made the restriction maps learnable. Schematically,
a layer is
\[
  X_{t+1}=X_t-\sigma\bigl(L_{\F(t)}\,(I\otimes W_1)\,X_t\,W_2\bigr),
\]
where the sheaf $\F(t)$ is produced from the features by a small network. By
Theorem~\ref{thm:hodge} the diffusion converges to the global sections of the learned sheaf, which
need not be constant. Bodnar et al.\ prove that suitable non-trivial sheaves separate classes of
nodes, in the diffusion limit, in settings where the constant sheaf cannot, and that richer families
of restriction maps separate more classes.

Three further references mark the range of the subject. \textcite{gebhart2023} show that standard
knowledge-graph embeddings are sheaves over the graph's schema, so that a consistent embedding is a
global section. \textcite{belfiore2021} place every feed-forward architecture in a Grothendieck
topos, with invariances as stacks; this is the most abstract and least tested proposal.
\textcite{ayzenberg2025} survey the whole area.

% -----------------------------------------------------------------------------------
\section{Reading guide}\label{sec:guide}

The order below assumes a reader who knows algebraic topology and homological algebra but has not
worked with sheaves on data.

\subsection{A first pass}
\begin{enumerate}
  \item \textcite{robinson2019}. A short expository article built around one worked problem with
        real measurements. It shows the dictionary of Section~\ref{sec:idea} in action.
  \item \textcite{curry2014}, the introduction and the chapters on cellular sheaves. This is the
        place to learn Definition~\ref{def:cellular}, its equivalence with constructible sheaves,
        and the dual notion of a cosheaf.
  \item \textcite{hansen2019}, for the Laplacian and Theorem~\ref{thm:hodge}. It is the technical
        basis of everything computational.
\end{enumerate}

\subsection{Three tracks}
\begin{description}[style=nextline, leftmargin=1.4em, itemsep=.5em]
  \item[Data fusion] \textcite{robinson2017} for the axiomatic argument; \textcite{robinson2020} for
        the consistency radius and its stability; \textcite{joslyn2020} for a complete application;
        \textcite{robinson2014} for the book-length account, including sampling and filtering.
  \item[Information theory] \textcite{abramsky2011}, then \textcite{vorobev1962} for the theorem
        behind it; \textcite{baudot2015} for the idea and \textcite{vigneaux2020} for the rigorous
        general framework; \textcite{peltre2020} for message passing.
  \item[Machine learning] \textcite{bodnar2022}, which is self-contained; \textcite{gebhart2023} for
        knowledge graphs; \textcite{ayzenberg2025} as the reference survey. \textcite{hansen2021}
        applies the same diffusion to opinion dynamics and is a good source of intuition.
\end{description}

\subsection{Background and history}
For the classical theory, \textcite{bredon1997} or the first chapters of \textcite{kashiwara1990};
for the categorical view, \textcite{maclane1992}. \textcite{rosiak2022} is an introduction written
around examples from outside pure mathematics, and \textcite{ghrist2014} places sheaves within
applied topology as a whole. For the history, \textcite{miller2000} on Leray and
\textcite{gray1979} for the period up to the 1970s.

\subsection{Software}
Two open libraries implement the constructions above: \texttt{pysheaf} (Python, by Robinson; sheaves
on finite spaces, consistency radius, and fusion by optimization) and \texttt{CellularSheaves.jl}
(Julia, in the AlgebraicJulia project; cellular sheaves and sheaf Laplacians).

\printbibliography[heading=bibintoc,title={References}]

\end{document}
